%% kapa: exam=1 round=37 year=2026 period=2 form=A subject=accounting subject_ko=회계학 seq=7 number=47
%% source: Q-Net 감정평가사 기출문제 (cst00309) artlSeq=5258924 file=raw/1008/5258924/2026년 제37회 감정평가사 1차 2교시.pdf
%% conversion: kordoc 4.24.51 -> tools/gen1.py (auto); review=proofread-sonnet; flags=none
\begin{question}{47}
(주)감평은 20×1년 4월 초 공장건물 신축공사를 시작하여 20×2년 9월 말 완공하였다. 동 건물은 차입원가 자본화 대상인 적격자산이다. 동 건물 신축 관련 자료가 다음과 같을 때, (주)감평이 20×1년에 자본화할 차입원가는? (단, 연평균지출액과 이자비용은 월할 계산하며 단수차이로 인한 오차가 있으면 가장 근사치를 선택한다.)
\begin{framedtext}
(1) 공장건물 신축을 위한 공사비 지출내역은 다음과 같다. \par
{\small\centering\begin{tabular}{@{}lr@{}}
일 자 & 금 액 \\\hline
20×1. 4. 1. & ￦200,000 \\
20×1. 10. 1. & 500,000 \\
20×2. 4. 1. & 500,000 \\
\end{tabular}\par}
(2) 차입금 현황은 다음과 같으며, 모든 차입금의 이자는 단리로 계산하여 만기일에 지급하는 조건이다. 차입금 A는 특정목적차입금이며, B, C는 일반목적차입금이다. \par
{\small\centering\begin{tabular}{@{}lrcr@{}}
구 분 & 차입금액 & 차입기간 & 연 이자율 \\\hline
차입금 A & ￦300,000 & 20×1. 4. 1. ∼ 20×2. 3. 31. & 8 \% \\
차입금 B & 240,000 & 20×1. 3. 1. ∼ 20×2. 2. 28. & 6 \% \\
차입금 C & 400,000 & 20×1. 7. 1. ∼ 20×2. 6. 30. & 5 \% \\
\end{tabular}\par}
\end{framedtext}
\begin{choices}
  \item ￦18,000
  \item ￦20,000
  \item ￦20,750
  \item ￦22,000
  \item ￦38,000
\end{choices}
\end{question}
\answer{3}
